In \Cref{cnstr:barconstruction}, we recalled that $\MO \langle 4n \rangle$ can be computed via a two-sided bar construction.  In this section we give a description of the bar construction, valid through a range of homotopy groups, which is particularly well-suited to the explicit identification of the $d_2$ in the bar spectral sequence as a Toda bracket.  Our main result is \Cref{thm:cofiber}.

\begin{comment}
In later sections, we combine \Cref{thm:cofiber} with the results of \Cref{sec:Goodwillie} to understand the unit map
$$\pi_{8n-1} \bS \to \pi_{8n-1} \MO \langle 4n \rangle,$$
and in particular to prove Theorem \ref{thm:intro-main}.
\end{comment}

\begin{cnstr} \label{cnstr:square}
Since $J$ is a map of non-unital rings in the homotopy category of spectra, the following diagram commutes up to homotopy:
$$
\begin{tikzcd}
\sOf^{\otimes 2} \arrow{r}{1 \otimes J} \arrow{d}{m} & \sOf \arrow{d}{J} \\
\sOf \arrow{r}{J} & \mathbb{S},
\end{tikzcd}
$$
where $m$ is the product map.
In the $\infty$-category of spectra, the fact that $J$ is a ring map is not a property, but actually additional structure.
In particular, there is a \emph{canonical homotopy} $a$ filling the above square:
$$
\begin{tikzcd}[column sep = huge]
\sOf^{\otimes 2} \arrow{dd}{m} \arrow{r}{1 \otimes J} & \sOf \arrow{dd}{J} \\ \\
\sOf \arrow{r}{J} \arrow[Leftrightarrow]{ruu}{a} & \bS.
\end{tikzcd}
$$
This homotopy $a$ may alternatively be viewed as a specific \emph{nullhomotopy} of $J \circ (1 \otimes J-m)$.
Let $P$ denote the homotopy cofiber of the map
$$
\begin{tikzcd}[column sep=huge]
\sOf^{\otimes 2} \arrow{r}{1 \otimes J-m} & \sOf
\end{tikzcd}
$$
Then the homotopy $a$ provides a canonical factorization
$$
\begin{tikzcd}[column sep=huge]
\sOf^{\otimes 2} \arrow{r}{1 \otimes J-m} & \sOf \arrow{d}{J} \arrow{r} & P \arrow[dashed]{dl} \\
& \mathbb{S}.
\end{tikzcd}
$$
\end{cnstr}

The main theorem of this section is the identification of the Thom spectrum $\MO \langle 4n \rangle$ with the cofiber of the above map $P \to \Ss$ in a range:

\begin{thm} \label{thm:cofiber}
Let $C$ denote the  cofiber of the map $P \to \mathbb{S}$ constructed above.
Then there is an equivalence of spectra $\tau_{\le 12n-2} C \simeq \tau_{\le 12n-2} \MO \langle 4n \rangle$.  Furthermore, the unit map
$$\tau_{\le 12n-2} \mathbb{S} \to \tau_{\le 12n-2} \MO \langle 4n \rangle$$ agrees with the natural map $\tau_{\le 12n-2} \mathbb{S} \to \tau_{\le 12n-2} C$.
\end{thm}

Before proving \Cref{thm:cofiber}, let us recall \Cref{cnstr:barconstruction}.  \Cref{cnstr:barconstruction} says that the spectrum $\MO \langle 4n \rangle$ can be calculated as the geometric realization of a two-sided bar construction
$$\MO \langle 4n \rangle = \left|\sBar(\bS,\spOf, \bS)_{\bullet} \right|.$$
Here, the action of $\spOf$ on the leftmost $\bS$ is via $\epsilon$, and the action on the rightmost $\bS$ is via $J_+$.

We may display this bar construction as a simplicial object,
$$
\begin{tikzcd}
\cdots \arrow[r,yshift=-3ex] \arrow[r,yshift=-1ex] \arrow[r,yshift=1ex] \arrow[r,yshift=3ex]
& \spOf^{\otimes 2} \arrow[r,yshift=-2ex,"\epsilon \otimes 1"] \arrow[r,"m"] \arrow[r,yshift=2ex,"1 \otimes J_+"]
& \spOf \arrow[r,yshift=-1ex,"\epsilon"] \arrow[r,yshift=1ex, "J_+"] 
& \bS,
\end{tikzcd}
$$
where the leftward degeneracy maps are omitted.  The key point is that, if we only wish to study $\tau_{\le 12n-2} \MO \langle 4n \rangle$, we need only study the partial simplicial diagram
$$
\begin{tikzcd}
\spOf^{\otimes 2} \arrow[r,yshift=-2ex,"\epsilon \otimes 1"] \arrow[r,"m"] \arrow[r,yshift=2ex,"1 \otimes J_+"]
& \spOf \arrow[r,yshift=-1ex,"\epsilon"] \arrow[r,yshift=1ex, "J_+"] 
& \bS.
\end{tikzcd}
$$
In the language of \cite[Lemma 1.2.4.17]{HA}, this is a diagram $\Delta^{\text{op}}_{\le 2} \to \Sp$.

\begin{lem} \label{lem:partialbar}
Let $X$ denote the  colimit of the partial simplicial diagram $\Delta^{\mathrm{op}}_{\le 2} \to \Sp$ given by $\mathrm{Bar}(\bS,\spOf,\bS)_{\le 2}$.  Then there is an equivalence of spectra $$\tau_{\le 12n-2} \MO \langle 4n \rangle \simeq \tau_{\le 12n-2} X.$$
\end{lem}

\begin{proof}
    Set $\mathrm{Bar}_{\leq k} = \abs{\sBar(\Ss, \spOf, \Ss)_{\leq k}}$. Then $\Ss = \sBar_{\leq 0}$ and $\MO \langle 4n \rangle$ is the  colimit of the $\sBar_{\leq k}$. We have a diagram
    \begin{center}
        \begin{tikzcd}[column sep=tiny]
            \Ss \arrow[r] & \sBar_{\leq 1} \arrow[r] \arrow[d] & \sBar_{\leq 2} \arrow[r] \arrow[d] & \sBar_{\leq 3} \arrow[r] \arrow[d] & \cdots \arrow[r] & \MO \langle 4n \rangle, \\
            & \Sigma \sOf & \Sigma^{2} \sOf^{\otimes 2} & \Sigma^{3} \sOf^{\otimes 3} & & 
        \end{tikzcd}
    \end{center}
    where the vertical maps are the  cofibers of the horizontal maps. The result now follows from the fact that, when $k \geq 3$, $\Sigma^{k} \sOf ^{\otimes k}$ is $(12n)$-connective.

%Applying $\pi_*$ to the full bar construction $\mathrm{Bar}(\bS,\spOf,\bS)$ yields a simplicial abelian group with layers
%$$\pi_*\left( \sBar(\bS,\spOf, \bS)_s \right) \cong \pi_*\left((\spOf)^{\otimes s} \right).$$
%The normalized chains of this simplicial abelian group form the $\sE^1$-page of the spectral sequence associated to the simplicial object $\sBar(\bS,\spOf, \bS)_\bullet$, which takes the form
%$$\sE^1_{s,t}=\pi_t\left((\sOf)^{\otimes s} \right) \Rightarrow \pi_{s+t} \MO \langle 4n \rangle.$$
%If $s \ge 3$, the lowest homotopy group of $\sOf^{\otimes s}$ is in dimension at least $12n-3$, and so nothing survives to the $\sE^2$-page with $s \ge 3$ and $s+t<12n$.  The classes with $s+t=12n$ can be sources of differentials, and so can contribute to the calculation of $\pi_{12n-1} \MO \langle 4n \rangle$.  However, no elements with $s \ge 3$ contribute to the calculation of $\pi_{\le 12n-2} \MO \langle 4n \rangle$.
\end{proof}

\begin{proof}[Proof of \Cref{thm:cofiber}]
In \Cref{lem:partialbar}, we established that $\tau_{\le 12n-2} \MO \langle 4n \rangle$ may be calculated as $\tau_{\le 12n-2}$ of the  colimit $X$ of a simplicial diagram
$$
\begin{tikzcd}
\spOf^{\otimes 2} \arrow[r,yshift=-2ex,"\epsilon \otimes 1"] \arrow[r,"m"] \arrow[r,yshift=2ex,"1 \otimes J_+"]
& \spOf \arrow[r,yshift=-1ex,"\epsilon"] \arrow[r,yshift=1ex, "J_+"] 
& \bS,
\end{tikzcd}
$$
where we have omitted the degeneracies from the notation.
%$$
%\begin{tikzcd}
%\spOf^{\otimes 2} \arrow[r,yshift=-2ex] \arrow[r] \arrow[r,yshift=2ex]
%& \spOf \arrow[r,yshift=-1ex] \arrow[l,yshift=1ex] \arrow[r,yshift=1ex] \arrow[l,yshift=-1ex]
%& \bS \arrow{l}.
%\end{tikzcd}
%$$
According to the cofinality statement of \cite[Lemma 1.2.4.17]{HA}, $X$ may be characterized as the lower right corner of the following cocartesian cube:
$$
\begin{tikzcd}[row sep={35,between origins}, column sep={45,between origins}]
      & \spOf^{\otimes 2} \ar[rr,"1 \otimes J_+"] \ar[dd,"\epsilon \otimes 1" near end] \ar[dl,"m"] & & \spOf \ar[dd,"\epsilon"] \ar[dl,"J_+"] \\
    \spOf \ar[rr,crossing over, "J_+" near end] \ar[dd,"\epsilon"] & & \mathbb{S}  \\
      & \spOf  \ar[rr, "J_+" near end] \ar[dl,"\epsilon"] & &  \mathbb{S}  \arrow{dl}  \\
    \mathbb{S} \arrow[rr,""] & & X. \ar[from=uu,crossing over]
\end{tikzcd}
$$
We finish the proof by showing that the $X$ appearing in the cube is equivalent to the spectrum $C$ from the theorem statement.  Indeed, taking fibers in the vertical direction, we learn that $X$ is the total cofiber of the square
$$
\begin{tikzcd}[column sep= huge]
\mathrm{fiber}(\spOf^{\otimes 2} \stackrel{\epsilon \otimes 1}{\longrightarrow} \spOf) \arrow{r} \arrow{d} &  \mathrm{fiber}(\spOf \stackrel{\epsilon}{\longrightarrow} \mathbb{S}) \arrow{d} \\
\mathrm{fiber}(\spOf \stackrel{\epsilon}{\longrightarrow} \mathbb{S}) \arrow{r} & \mathbb{S},
\end{tikzcd}
$$
which simplifies to the square
$$
\begin{tikzcd}[column sep= huge]
\sOf \oplus \sOf^{\otimes 2} \arrow{r}{(1,1 \otimes J)} \arrow{d}{(1,m)} &  \sOf \arrow{d}{J} \\
\sOf \arrow{r}{J} \arrow[Leftrightarrow]{ru}{1\oplus a} & \mathbb{S}.
\end{tikzcd}
$$
The pushout of the arrows $(1,1 \otimes J)$ and $(1,m)$ is calculated as the cofiber of the map
$$
\begin{tikzcd} [column sep = huge, ampersand replacement=\&]
\sOf \oplus \sOf^{\otimes 2} \arrow{r}{\begin{pmatrix} 1 & 1 \otimes J \\ -1 & -m \end{pmatrix}} \& \sOf \oplus \sOf,
\end{tikzcd}
$$
or equivalently the cofiber of the map
$$
\begin{tikzcd} [column sep = huge]
\sOf^{\otimes 2} \arrow{r}{1 \otimes J-m} & \sOf.
\end{tikzcd}  
$$
This cofiber is the spectrum $P$ of the theorem statement.  To obtain the final sentence of the theorem, note that the unit map from $\bS$ to $\MO \langle 4n \rangle$ is the map from $\mathrm{Bar}(\bS,\spOf,\bS)_0$ into the geometric realization of the full bar construction.  This factors through the partial bar construction $|\mathrm{Bar}(\bS,\spOf,\bS)_{\le 2}| \simeq X$, via the map $\bS \to X$ that appears three times in the above cocartesian cube.
\qedhere
\end{proof}

%%% Local Variables:
%%% mode: latex
%%% TeX-master: "Main"
%%% eval: (visual-line-mode t)
%%% eval: (auto-fill-mode 0)
%%% End:
